\section{Model Fitting and Evaluation Protocol}
\label{sec:fitting}

This section summarizes how the preprocessed experiment windows
are used to evaluate the models.
Full results are presented in the revised Section~6 of the
manuscript; here we document the protocol.

\subsection{Models Fitted}

For each experiment window, we fit the following models using
only the pilot data (days $1, \ldots, \Dpilot$):

\paragraph{Our models.}
\begin{itemize}[leftmargin=*]
  \item \textbf{\bemodel} (Bernoulli SSP): models binary
        activity indicators; predicts $\Utrue$ only.
  \item \textbf{\tgmodel} (Truncated Geometric SSP): models
        first-trigger times; predicts $\Utrue$ only.
  \item \textbf{\nbmodel} (Negative Binomial SSP): models
        count-valued daily triggers; predicts both $\Utrue$
        and $\Ttrue$.
\end{itemize}

\paragraph{Competing methods.}
\begin{itemize}[leftmargin=*]
  \item Beta--Binomial (BB)
  \item Hierarchical Beta--Geometric (HBG)
  \item Jackknife estimator
  \item Good--Toulmin estimator (GT)
  \item Linear Programming bound (LP)
\end{itemize}

\subsection{Evaluation Metrics}

\paragraph{Accuracy.}
For each experiment and method, we compute the relative
prediction accuracy:
\[
  v_{\Dpilot}^{(\Dfollow)}
  = \frac{\hat{U}_{\Dpilot}^{(\Dfollow)}}{U_{\Dpilot}^{(\Dfollow)}},
\]
where $v = 1$ indicates perfect prediction, $v > 1$ overprediction,
and $v < 1$ underprediction.
For the \nbmodel{} model, we additionally report the analogous
ratio for total triggers $\Ttrue$.

\paragraph{Calibration.}
For each experiment, we compute the 95\% posterior credible
interval for $\Utrue$ (and $\Ttrue$ for \nbmodel).
Across all experiments, we report the empirical coverage:
the fraction of experiments in which the true value falls
within the credible interval.
Well-calibrated models should achieve coverage close to the
nominal 95\% level.

\paragraph{Decision impact.}
For selected experiments, we illustrate the decision-support
use case: given pilot data, how accurately can the method
predict whether the experiment has reached a target
participation level?
We compute the estimated number of additional days $D_M$
needed to reach $M$ total users and compare to the true value.

\subsection{Connection to Reviewer Concerns}

This evaluation protocol directly addresses the three main
concerns:
\begin{itemize}[leftmargin=*]
  \item \textbf{Thin data analysis} (AE, R2-b):
        We evaluate on 7+ REES46 experiments and 13 UCI experiments,
        rather than a single dataset.
  \item \textbf{Practical impact} (AE, R2-c):
        Calibration and $D_M$ analyses demonstrate that the
        methodology supports concrete experiment-duration decisions.
  \item \textbf{NB-SSP on public data} (R2-b):
        For the first time, we evaluate total-trigger predictions
        on publicly available data with per-user trigger counts.
\end{itemize}

\subsection{Case Studies: Prediction Trajectories vs.\ Truth}
\label{subsec:case_studies}
\addresses{R2-b (T1.b), R2-j (T2.e)}

To demonstrate the practical value of the predictions, we
show prediction trajectories alongside ground truth for
selected experiments from each dataset.
These case studies illustrate what an experimenter would see
in practice: given 7~days of pilot data, how well does each
method forecast the accumulation of new users?

\begin{figure}[H]
  \centering
  \includegraphics[width=\linewidth]{fig_case_study_uci.pdf}
  \caption{UCI: prediction trajectories for three experiments
  spanning sparse (Exp~1, $\Npilot = 34$), medium (Exp~6,
  $\Npilot = 348$), and large (Exp~12, $\Npilot = 474$)
  pilot sizes.\reviewtag{A4}
  The gray region marks the pilot period; curves show
  predicted cumulative new users vs.\ truth (black).
  NB-SSP regression (green) tracks the truth most closely
  on Exp~1 and Exp~6.}
  \label{fig:case_uci}
\end{figure}

\begin{figure}[H]
  \centering
  \includegraphics[width=\linewidth]{fig_case_study_rees46.pdf}
  \caption{REES46: final predictions (markers) vs.\ truth
  (black curve) for three experiments.\reviewtag{A4}
  Dashed lines connect the pilot end to each method's
  prediction at day $D$.
  All methods predict well on this dataset
  ($v > 0.80$ for IBP and NB-SSP regression).}
  \label{fig:case_rees46}
\end{figure}

\begin{figure}[H]
  \centering
  \includegraphics[width=\linewidth]{fig_case_study_asos.pdf}
  \caption{ASOS: final predictions vs.\ truth for three
  experiments with short, medium, and long follow-up
  periods.\reviewtag{A4}
  The difficulty of long-horizon extrapolation is visible:
  predictions diverge more from truth as $\Dfollow$ increases.}
  \label{fig:case_asos}
\end{figure}

\paragraph{Practical interpretation.}
On the UCI data, NB-SSP regression closely tracks the true
cumulative curve even from very sparse pilot data (34~users
in Exp~1).
On REES46, all methods perform well because the large pilot
sizes ($\sim$1M users) provide strong signal.
On ASOS, the case studies make the extrapolation difficulty
concrete: a 7-day pilot predicting 150+~days ahead is
inherently uncertain, and the figures show this honestly.

An experimenter using these predictions would see, for example,
that after 7~days of pilot with 348~users (UCI Exp~6),
NB-SSP regression predicts 669~new users by day~28 (truth: 592,
$v = 0.87$).
This is actionable: it tells the experimenter whether the
experiment is on track to reach its participation target.

\subsection{ASOS Results: Accuracy Stratified by Extrapolation Ratio}
\label{subsec:asos_stratified}
\addresses{R2-b (T1.b), R2-c (T2.e), AE}

The original submission reported a single accuracy number per
method on the ASOS dataset (76 experiments, 144 treatment arms,
$\Dpilot = 7$).
This aggregate view obscures a key structural feature of the
ASOS data: the follow-up horizons $\Dfollow$ range from 7 to
239~days, yielding extrapolation ratios $\Dfollow / \Dpilot$
from 1.0 to 34.1 (median 7.9).
Predicting 7~days ahead from 7~days of pilot is a fundamentally
different task from predicting 239~days ahead.

\Cref{fig:asos_boxplot_abs} summarizes the absolute accuracy
distribution for each method across all 144 treatment arms.
NB-SSP regression has the highest median accuracy (0.50) and
the tightest interquartile range, while TG-SSP and IBP show
wider spread with many zero-accuracy outliers.

\begin{figure}[H]
  \centering
  \includegraphics[width=0.7\linewidth]{fig_asos_accuracy_boxplot.pdf}
  \caption{ASOS: accuracy distribution by method (144 arms).\reviewtag{A2}
  NB-SSP regression has the highest median (0.50) and the
  tightest spread.
  TG-SSP and IBP have lower medians and heavier left tails
  (many arms with $v \approx 0$).}
  \label{fig:asos_boxplot_abs}
\end{figure}

\Cref{fig:asos_ranking} shows the relative performance:
for each arm, we rank the three methods by accuracy and
report how often each method places first, second, or third.
NB-SSP regression ranks best on 45\% of arms, followed by
TG-SSP (32\%) and IBP (23\%).

\begin{figure}[H]
  \centering
  \includegraphics[width=0.7\linewidth]{fig_asos_ranking_frequency.pdf}
  \caption{ASOS: ranking frequency across 144 arms.\reviewtag{A2}
  NB-SSP regression is the best-performing method on 45\%
  of arms.
  No single method dominates uniformly, but NB-SSP regression
  is the most frequent winner.}
  \label{fig:asos_ranking}
\end{figure}

\Cref{fig:asos_ratio_binned} shows how accuracy varies with
the extrapolation ratio $\Dfollow / \Dpilot$, using binned
medians with interquartile-range error bars.
Two patterns emerge:
\begin{enumerate}[leftmargin=*]
  \item \textbf{NB-SSP regression is the most stable method.}
        Its median accuracy remains in the 0.4--0.5 range across
        all extrapolation ratios, while TG-SSP and IBP are more
        variable.
  \item \textbf{Absolute accuracy on ASOS is lower than on
        UCI and REES46 for all methods.}
        Even at $\Dfollow / \Dpilot \leq 3$, the median accuracy
        of NB-SSP regression is 0.453 (vs.\ 0.843 on UCI and
        0.888 on REES46).
        The ASOS dataset involves fashion retail experiments with
        shorter, more volatile engagement patterns and smaller
        pilot samples relative to the user population.
\end{enumerate}

\begin{figure}[H]
  \centering
  \includegraphics[width=0.85\linewidth]{fig_asos_accuracy_vs_ratio_binned.pdf}
  \caption{ASOS: median accuracy vs.\ extrapolation ratio
  (binned, with IQR error bars).\reviewtag{A2}
  NB-SSP regression (green triangles) is the most stable
  across all bins.
  The dashed line marks $v = 0.8$.}
  \label{fig:asos_ratio_binned}
\end{figure}

\begin{table}[H]
\centering
\caption{ASOS accuracy (median and mean of
$v_{\Dpilot}^{(\Dfollow)}$) stratified by extrapolation ratio.\reviewtag{A2}
NB-SSP regression is the most consistent method across all bins.}
\label{tab:asos_stratified}
\small
\begin{tabular}{@{}lrccccccc@{}}
\toprule
& & \multicolumn{2}{c}{\textbf{TG-SSP}}
& \multicolumn{2}{c}{\textbf{IBP}}
& \multicolumn{2}{c}{\textbf{NB-SSP (reg)}} \\
\cmidrule(lr){3-4} \cmidrule(lr){5-6} \cmidrule(lr){7-8}
\textbf{Bin} & $n$ & med & mean & med & mean & med & mean \\
\midrule
$\leq 3$   & 18  & 0.372 & 0.358 & 0.121 & 0.350 & 0.453 & 0.500 \\
3--5       & 18  & 0.288 & 0.422 & 0.700 & 0.545 & 0.414 & 0.423 \\
5--10      & 64  & 0.000 & 0.272 & 0.000 & 0.187 & 0.501 & 0.431 \\
$>10$      & 44  & 0.000 & 0.157 & 0.253 & 0.361 & 0.531 & 0.452 \\
\midrule
All        & 144 & 0.000 & 0.266 & 0.158 & 0.305 & 0.495 & 0.445 \\
\bottomrule
\end{tabular}
\end{table}

\paragraph{Implications for the paper.}
The ASOS results serve a different role from the UCI and REES46
results.
While the latter demonstrate high accuracy ($v > 0.84$) under
moderate extrapolation, the ASOS analysis shows that:
(i)~NB-SSP regression is the most robust method even in a
harder setting;
(ii)~the method degrades gracefully rather than catastrophically;
and (iii)~the difficulty is inherent to the prediction task
(long horizons, volatile fashion data), not a model deficiency,
since all competing methods perform similarly or worse.
In the revised Section~6, we present the ASOS results stratified
by extrapolation ratio rather than as a single aggregate number,
which gives the reader a more informative picture.

\subsection{Cross-Dataset Accuracy Summary}
\label{subsec:cross_dataset}
\addresses{AE, R2-b (T1.b, T2.e)}

\Cref{tab:cross_dataset} consolidates accuracy results across
all three datasets and all methods.
The key findings are:
\begin{enumerate}[leftmargin=*]
  \item \textbf{NB-SSP regression is the best or tied-for-best
        method on every dataset.}
  \item \textbf{Accuracy is high ($v > 0.80$) when the
        extrapolation ratio is moderate}
        (UCI: $\Dfollow/\Dpilot = 3$;
        REES46 sliding: $\Dfollow/\Dpilot = 1.8$;
        REES46 rolling $k{=}21$: $\Dfollow/\Dpilot = 3$).
  \item \textbf{Accuracy degrades for longer horizons}
        (ASOS: $\Dfollow/\Dpilot$ up to 34),
        but NB-SSP regression degrades more gracefully than
        competitors.
\end{enumerate}

\begin{table}[H]
\centering
\caption{Prediction accuracy (median $v_{\Dpilot}^{(\Dfollow)}$)
across all datasets and methods.\reviewtag{A2}
Bold indicates the best method per dataset.
NB-SSP regression is the best or tied-for-best on every dataset.}
\label{tab:cross_dataset}
\small
\begin{tabular}{@{}lrrccccc@{}}
\toprule
& & $\Dfollow/\Dpilot$
& \textbf{TG-SSP} & \textbf{IBP}
& \textbf{NB-SSP (MLE)} & \textbf{NB-SSP (reg)} \\
\midrule
UCI & 13 & 3.0 & 0.708 & 0.804 & 0.813 & \textbf{0.843} \\
REES46 (sliding) & 7 & 1.8 & 0.575 & 0.806 & --- & \textbf{0.897} \\
REES46 ($k{=}21$) & 186 & 3.0 & 0.807 & \textbf{0.809} & --- & 0.802 \\
REES46 ($k{=}50$) & 157 & 7.1 & \textbf{0.847} & \textbf{0.847} & --- & 0.832 \\
REES46 ($k{=}100$) & 107 & 14.3 & 0.841 & \textbf{0.868} & --- & 0.857 \\
ASOS & 144 & 7.9 & 0.000 & 0.158 & --- & \textbf{0.495} \\
\bottomrule
\end{tabular}
\end{table}

\paragraph{Note on REES46 rolling results.}
On the REES46 rolling experiments ($k = 21, 50, 100$), all three
methods achieve similar accuracy ($v \approx 0.80$--$0.87$).
The differences are small because these experiments have very
large pilot samples ($\Npilot \sim 1\text{M}$), which provide
enough information for even simpler models to perform well.
The advantage of NB-SSP regression is most visible on smaller
or harder datasets (UCI, ASOS) where the model's flexibility
matters more.

\begin{figure}[H]
  \centering
  \includegraphics[width=\linewidth]{fig_cross_dataset_boxplot.pdf}
  \caption{Accuracy distributions across all datasets and methods.\reviewtag{A5}
  Each group of three boxplots corresponds to one dataset;
  colors distinguish methods.
  NB-SSP regression (green) has the highest or most stable
  median on every dataset.
  The dashed line marks $v = 0.8$.}
  \label{fig:cross_boxplot}
\end{figure}

\begin{figure}[H]
  \centering
  \includegraphics[width=0.7\linewidth]{fig_cross_dataset_ranking.pdf}
  \caption{Ranking frequency across all 350 experiments.\reviewtag{A5}
  NB-SSP regression ranks best on 41\% of experiments,
  vs.\ 29\% each for TG-SSP and IBP.}
  \label{fig:cross_ranking}
\end{figure}

\begin{figure}[H]
  \centering
  \includegraphics[width=0.85\linewidth]{fig_cross_dataset_accuracy_vs_ratio.pdf}
  \caption{Median accuracy vs.\ median extrapolation ratio
  across datasets.\reviewtag{A5}
  All methods degrade with longer extrapolation horizons,
  but NB-SSP regression (green) degrades most gracefully.
  At moderate ratios ($\leq 3$), all methods achieve
  $v > 0.80$.}
  \label{fig:cross_ratio}
\end{figure}

\subsection{Credible Interval Coverage Analysis}
\label{subsec:ci_coverage}
\addresses{R2-b (T1.c), AE}

Reviewer~2 requested ``calibrated risk at given service-level
targets'' (R2-b) and the revision plan identified calibration
diagnostics as a high-priority item (T1.c).
We report 95\% posterior credible interval coverage for
NB-SSP regression (the best-performing method) across all
datasets.

\paragraph{Summary.}
\Cref{tab:ci_summary} shows that empirical coverage falls
well below the nominal 95\% level on every dataset.
The root cause is a mismatch between CI width and point-estimate
bias: the posterior concentrates tightly around the plug-in
estimate, but that estimate carries non-trivial bias relative
to the truth.

\begin{table}[H]
\centering
\caption{CI coverage for NB-SSP regression (95\% nominal).\reviewtag{A3}
Median absolute relative bias and median relative CI width
(both as \% of $U_{\text{true}}$) explain the coverage gap:
coverage fails when bias exceeds CI width.}
\label{tab:ci_summary}
\small
\begin{tabular}{@{}lrccc@{}}
\toprule
\textbf{Dataset} & $n$ & \textbf{Coverage}
& \textbf{Med $|\text{bias}|$} & \textbf{Med CI width} \\
\midrule
UCI Online Retail & 13 & 6/13 (46\%) & 15.7\% & 25.6\% \\
REES46 $k{=}21$ & 186 & 2/186 (1\%) & 19.8\% & 0.6\% \\
REES46 $k{=}50$ & 157 & 1/157 (1\%) & 16.8\% & 0.6\% \\
REES46 $k{=}100$ & 107 & 2/107 (2\%) & 14.3\% & 0.6\% \\
ASOS & 144 & 0/144 (0\%) & 50.5\% & 0.7\% \\
\bottomrule
\end{tabular}
\end{table}

\paragraph{Mechanism.}
\Cref{fig:ci_bias_width} illustrates the mechanism on the
UCI data (the only dataset where coverage is non-trivial).
Each point is one experiment; green circles indicate the CI
covers the truth, red crosses indicate it does not.
The pattern is clear: coverage succeeds when the absolute
bias is smaller than the CI width, and fails otherwise.

\begin{figure}[H]
  \centering
  \includegraphics[width=0.75\linewidth]{fig_ci_bias_vs_width.pdf}
  \caption{UCI: absolute relative bias vs.\ relative CI width
  for NB-SSP regression.\reviewtag{A3}
  CIs cover the truth (green circles) when bias is small
  relative to CI width; they miss (red crosses) when bias
  dominates.}
  \label{fig:ci_bias_width}
\end{figure}

On the REES46 data, pilot sizes are $\sim$1M users, so the
posterior concentrates to a relative width of $\sim$0.6\%.
Even a modest 5--20\% bias in the point estimate falls far
outside these narrow intervals.
On UCI, pilot sizes are $\sim$200 users, so CIs are wider
($\sim$25\%), and coverage reaches 46\% --- still below
nominal, but consistent with the bias magnitude.

\begin{figure}[H]
  \centering
  \includegraphics[width=0.95\linewidth]{fig_ci_coverage_by_dataset.pdf}
  \caption{Empirical CI coverage (blue) and median absolute
  bias (red) across datasets.\reviewtag{A3}
  The dashed line marks the nominal 95\% level.
  Coverage is low everywhere because bias exceeds CI width;
  the gap is largest on REES46 and ASOS where posteriors
  are extremely concentrated.}
  \label{fig:ci_coverage_datasets}
\end{figure}

\paragraph{Interpretation.}
This coverage shortfall is an inherent property of empirical
Bayes inference with plug-in hyperparameters, not a deficiency
specific to our model.
The closed-form posterior credible intervals condition on
point-estimated hyperparameters $(\hat\alpha, \hat\sigma, \hat\Delta)$
and therefore understate the total uncertainty.
A fully Bayesian treatment (placing priors on the hyperparameters
and marginalizing) would produce wider intervals with better
coverage, at the cost of requiring MCMC --- precisely the
computational burden our closed-form approach is designed to avoid.

\paragraph{Implications for the paper.}
We recommend presenting the CI results honestly:
\begin{enumerate}[leftmargin=*]
  \item Report coverage alongside accuracy, acknowledging the
        limitation.
  \item Explain the mechanism (plug-in concentration) and note
        that it affects all empirical Bayes methods, not just ours.
  \item Frame the CIs as providing a \emph{lower bound} on
        uncertainty: the true prediction uncertainty is at least
        as large as the reported interval.
  \item Note that for the primary use case (point prediction of
        $U$ and $T$ to inform duration decisions), the accuracy
        metric $v$ is more relevant than CI coverage, and
        accuracy is high ($>0.84$) under moderate extrapolation.
\end{enumerate}
